% =====================================================================
% LE COURS, VERSION ENSAE : parties I a IX, recale le 19 octobre 2026 sur le rapport ENSAE de 100 pages.
% Le preambule est commun avec cours_memoire_questions_ensae.tex, voir
% cours_preambule.tex : ne pas le dupliquer.
% =====================================================================
\input{cours_preambule}

\title{\vspace{-1.3cm}\color{navy}\bfseries Le mémoire (version ENSAE), concept par concept\\[3pt]
  \large Ce que chaque outil résout, comment il marche,
  et la phrase qui l'explique en trente secondes}
\author{Kélian Kaddouri}
\date{Mise à jour du 28 octobre 2026}

\begin{document}
\maketitle
\vspace{-0.8cm}

\begin{cle}
\textbf{À quoi sert ce document.} Le mémoire empile une vingtaine d'outils, chacun choisi
pour une raison précise, et il est facile de les manipuler correctement sans savoir les
\emph{défendre}. Chaque section suit le même gabarit~: le problème que l'outil résout,
l'idée en langage courant, la formule, ce que ça donne dans notre cas, la phrase à dire à
voix haute quand on vous interroge, et le piège à éviter.

\textbf{La phrase encadrée en vert est le livrable de chaque section.} Si vous ne retenez
qu'elle, vous tenez l'oral.

\textbf{Un second document, \texttt{cours\_memoire\_questions\_ensae}, reprend le rapport ENSAE chapitre par chapitre (partie~X).} Pour chacun : le chapitre en une
phrase, un encadré de ce qu'il faut savoir \emph{par c\oe ur}, puis les questions que le jury peut
poser, avec la réponse à tenir et le folio où s'appuyer. C'est le document à relire la veille de la
soutenance.
\end{cle}

\smallskip
\noindent\textcolor{accent}{\textbf{Base ENSAE du 19 octobre ; réponses actualisées au 28 octobre 2026.}}
Ce cours conserve les folios du rapport ENSAE de $100$~pages (\texttt{memoire/main\_ensae.pdf}),
corrigé le 19 octobre. Les réponses intègrent les rectifications de portée les plus récentes.
Une partie finale, explicitement hors rapport déposé, présente les compléments Institut du 27 octobre.
L'exposé ENSAE doit tenir en \textbf{15 minutes au maximum}, questions séparées. Trois points ont changé depuis la version pour le manuscrit de $199$~pages :
la \textbf{valeur de l'information} (la réduction de la plage par un registre d'incidents est de
$32{,}7\,\%$ en moyenne, de $22$ à $44\,\%$ selon la graine, et non $66\,\%$ ; aucune paire ne se
détache, quatre dépendances levant chacune de $11{,}2$ à $13{,}3\,\%$) ; le \textbf{transfert du
risque} (coût de financement de $8{,}0$ à $2{,}9$~M\euro, soit $-64\,\%$ ; pertes conservées
comprises, de $11{,}67$ à $5{,}12$, soit $-56\,\%$ ; seuil de $64\,\%$ de portée sans allègement de
capital, et non $78\,\%$) ; et la \textbf{part insensible au sens} ($79\,\%$ du capital : socle
$61\,\%$ et plancher commun aux sommets $18\,\%$). Rappel des erreurs de fond déjà corrigées à
l'oral : le capital est une $\mathrm{VaR}_{99{,}5\,\%}$ \emph{non centrée} ; seul le \emph{groupe de
tête}, P1 et P4, résiste à l'ignorance de la direction ; l'écart entre les deux états de conformité
est \emph{posé} pour les quatre canaux, seuls les niveaux de référence de la fréquence et de la
détection étant calibrés. Les nombres du cours se retrouvent dans le rapport, ou dans le supplément
en ligne, ou dans les sorties versionnées.

\bigskip
\noindent\textbf{Le schéma d'ensemble.} Cinq étages, de la donnée au capital. Chaque partie de ce
document en traite un.

\begin{center}
\begin{tikzpicture}[node distance=4.5mm]
\node[brique] (d) {\textbf{Données}\\OpRisk, PRC\\Hackmageddon};
\node[brique,right=of d] (f) {\textbf{Fréquence}\\binomiale\\négative};
\node[brique,right=of f] (s) {\textbf{Sévérité}\\seuil\,+\,excès GPD\\$\xi=0{,}5954$};
\node[brique,right=of s] (c) {\textbf{Cascade}\\dirigée $W$\\$\rho(W)<1$};
\node[brique,right=of c] (m) {\textbf{Conformité}\\multi-états\\(Markov)};
\node[brique,right=of m,fill=vert!15,draw=vert!70!black] (k)
  {\textbf{Capital}\\$\mathrm{VaR}_{99{,}5\%}(L)$};
\draw[fleche] (d)--(f); \draw[fleche] (f)--(s); \draw[fleche] (s)--(c);
\draw[fleche] (c)--(m); \draw[fleche] (m)--(k);
\draw[flechef] (d.south) to[out=-60,in=-120] node[etiq,below,pos=0.5]
  {la direction de $W$ n'est PAS identifiable ici} (c.south);
\end{tikzpicture}
\end{center}

\noindent La flèche rouge situe l'information manquante sur les liens entre piliers.
Les lois de pertes et la fréquence de référence sont ajustées sur les données ; la force des
liens et la correspondance entre conformité et paramètres restent posées. L'identification
partielle examine ici la direction, conditionnellement à la force symétrique retenue.

\tableofcontents
\newpage

% =====================================================================
\part*{I. D'où viennent les nombres}
\addcontentsline{toc}{part}{I. D'où viennent les nombres}

\section{Les trois sources, et ce que chacune peut dire}

\lequoi{Pour chiffrer un capital, il faut trois choses~: combien d'incidents par an, combien
coûte un incident, et comment les incidents s'enchaînent. Aucune source publique ne porte
les trois.}

\lidee{On ne cherche pas \emph{la} bonne base, on répartit les questions entre les sources
selon ce que chacune peut porter, et on déclare ce qui manque.}

\cheznous{Trois bases pour les marges, un corpus pour l'ordre, deux sources de contrôle.
\begin{itemize}[itemsep=1pt,topsep=2pt]
\item \textbf{SAS OpRisk Global Data}~: pertes opérationnelles en dollars, catégories de Bâle ;
$582$ pertes cyber $\times$ finance sur 2000--2026 pour la sévérité, $4\,329$ transitions
annuelles pour le test de direction. Elle porte le \textbf{niveau} de sévérité et la fréquence
d'entrée. Son pas est annuel, ce qui efface l'ordre des événements.
\item \textbf{Privacy Rights Clearinghouse (PRC)}~: $15\,053$ incidents retenus, exprimés en
\emph{nombre d'enregistrements compromis}, pas en euros. Elle sert d'\textbf{axe de sensibilité}
après conversion, et de clé de répartition ($\lambda_{ref}=341$).
\item \textbf{Hackmageddon}~: $1\,041$ incidents, $840$ à vecteur identifié. Elle donne les
\textbf{parts par vecteur d'attaque}, jamais un niveau ; c'est une citation, non recalculable.
\item \textbf{Dix rapports post-incident} ($22$ transitions)~: peu de volume, mais une enquête y
reconstitue la \emph{séquence} des défaillances. C'est la seule source qui porte
l'\textbf{ordre}.
\item \textbf{VERIS} ($10\,591$ incidents) montre que le bon schéma existe sans être rempli, et
\textbf{quatre SFCR} donnent l'échelle d'entité.
\end{itemize}}

\laphrase{Aucune base ne dit tout. OpRisk donne le niveau mais pas l'ordre, PRC donne du volume
mais pas d'euros, Hackmageddon la structure par vecteur, et les rapports d'enquête donnent l'ordre
mais pas le volume. Je les utilise pour ce que chacune peut porter, et je dis explicitement ce qui
reste absent.}

\lepiege{Ne jamais laisser croire qu'on a \og une base de données DORA\fg{}. Elle n'existe pas.
Le règlement met en place le reporting qui la produira, et c'est justement l'objet de la
recommandation finale du mémoire.}

\section{La conversion Jacobs}

\lequoi{PRC compte des enregistrements compromis, pas des euros. Il faut passer de l'un à
l'autre.}

\lidee{Une régression publiée établit, sur un échantillon d'incidents dont on connaît à la
fois le nombre d'enregistrements et le coût, une relation log-log entre les deux.}

\laformule{$\ln L = 7{,}68 + 0{,}76\,\ln X + \varepsilon$, où $X$ est le nombre
d'enregistrements et $L$ la perte en dollars. Le résidu a un écart-type de $0{,}523$ en logarithme :
à volume donné, le coût réel varie d'un facteur $2{,}4$ autour de la droite, à $90\,\%$.}

\cheznous{L'exposant $0{,}76 < 1$ dit que le coût croît \emph{moins vite} que le nombre
d'enregistrements~: un incident dix fois plus gros ne coûte pas dix fois plus cher, mais
environ six fois. Il y a un coût fixe important dans tout incident.}

\laphrase{Jacobs n'est pas un lien fréquence-sévérité, c'est une conversion de sévérité
seule~: elle transforme un nombre d'enregistrements en montant. Et mon capital principal
n'en dépend pas, il est calculé sur OpRisk, qui est déjà en dollars.}

\lepiege{Le modèle traitait d'abord cette conversion comme \emph{exacte}, en ignorant le résidu.
En le réinjectant enregistrement par enregistrement, le niveau du surcoût PRC bouge de $7\,\%$ mais
l'intervalle ne s'élargit presque pas, le bruit se moyennant dans l'ajustement GPD ($\xi$ passe de
$1{,}03$ à $1{,}07$, matériel complémentaire). Et ne jamais lire $\xi_{\text{PRC}}=1{,}03$ comme un
risque plus lourd : une puissance multiplie l'indice de queue par son exposant, donc cet indice
décrit la queue des \emph{enregistrements} transportée en euros.}

% =====================================================================
\section{Le seuil de collecte, et les trois fréquences}

\lequoi{Trois nombres du mémoire portent le nom de \og{}fréquence\fg{} et vont de $0{,}092$ à
$341$ par an. Confondre deux d'entre eux fait dire n'importe quoi.}

\lidee{Ils ne comptent pas la même chose, parce qu'ils ne sont pas lus au même endroit. Une
base ne retient que les pertes au-dessus d'un \textbf{seuil de collecte} : le taux qu'on y lit
est celui des incidents \emph{matériels}, générateurs d'une perte enregistrée, et non celui de
tous les incidents au sens du reporting réglementaire.}

\cheznous{\;\\[-6pt]
\begin{center}\small
\begin{tabular}{@{}l l l@{}}
\toprule
$\lambda_{ref}=341$/an & notifications PRC, organisations financières & \emph{clé de répartition} par vecteur \\
$\lambda=21{,}6$/an & incidents TIC matériels, secteur & \textbf{l'entrée du moteur} \\
$\lambda_{\text{entité}}=0{,}092$/an & les mêmes, lus à la taille de l'entité & l'échelle d'une entité \\
\bottomrule
\end{tabular}
\end{center}
Les seaux du panel valent $0{,}10$ (toutes firmes) et $0{,}21$ (firmes à au moins dix événements),
mais ils ne décrivent aucune entité : la lecture juste estime $\lambda$ en fonction de la taille et
la lit à celle de l'entité visée. La base PRC entière compte $2\,150$ notifications par an ; le
$341$ ne porte que sur les organisations financières.}

\laphrase{Le facteur $3\,719$ entre ma plus grande et ma plus petite fréquence n'est pas une
incohérence, c'est trois niveaux d'observation. Un seul alimente le moteur, les deux autres servent
à répartir et à interpréter. Le seuil de collecte est ce qui les sépare.}

\lepiege{Le taux $\lambda$ est une propriété d'\textbf{entité}, lisible sur un panel
firme-année ; la charge systémique commune, elle, ne se lit qu'au niveau \textbf{agrégé} et
après retrait de tendance. Les confondre revient à imputer au risque systémique ce qui n'est
qu'un effet de taille des firmes.}

\section{Mesurer un biais au lieu de le déclarer}

\lequoi{Tout le monde déclare que sa base est biaisée. Presque personne ne dit de combien. Or
déclarer ne coûte rien, mesurer engage.}

\lidee{Deux biais classiques d'une base de pertes se mesurent avec la base elle-même. Le
\textbf{biais de taille}, parce que la base porte les actifs de la firme sinistrée. La
\textbf{troncature à droite}, parce qu'elle porte la date de saisie en plus de l'année de
survenance.}

\laformule{Pour la taille, une régression log-log $\ln L = a + b\ln(\text{actifs})$ donne
l'élasticité $b$. Mais elle est biaisée par le seuil de collecte : une petite firme n'entre
dans la base que si sa perte est anormalement grande, ce qui \emph{aplatit} la pente. On
corrige par un maximum de vraisemblance sur une lognormale \textbf{tronquée à gauche} en $u_0$,
où la contribution de chaque observation est divisée par $1-\Phi\bigl((\ln u_0-\mu)/s\bigr)$.
Pour la troncature à droite, on construit le triangle survenance $\times$ délai de remontée et
l'on applique un \textbf{chain-ladder} aux comptes.}

\cheznous{Élasticité brute $0{,}042$ avec un $R^2$ de $0{,}005$ (matériel complémentaire), ce qui
\emph{réfuterait} l'explication usuelle. Corrigée de la troncature, elle double à
$\mathbf{0{,}087}$ $[0{,}026\,;0{,}148]$ et devient significative, mais reste faible : la sévérité
transposée à l'entité notionnelle n'est multipliée que par $0{,}85$ environ ($0{,}854$ au
chapitre~10). Côté remontée, le délai médian est de trois ans (matériel complémentaire) et les
$97$ événements observés sur 2023--2025 deviennent \textbf{$255$} à l'ultime, soit $+163\,\%$.}

\laphrase{Mes deux biais les plus souvent opposés, je ne me contente pas de les signaler : je
les chiffre. Le biais de taille coûte quinze pour cent environ, pas un facteur dix. Et la baisse
apparente des incidents après 2022 n'est pas une baisse du risque, c'est une remontée
incomplète que le chain-ladder rattrape.}

\lepiege{Les deux corrections vont en sens opposés et se compensent presque exactement sur la
\emph{perte moyenne} ($0{,}861\times1{,}147=0{,}987$). Elles ne se compensent \textbf{pas} sur
le capital : à queue lourde, le quantile suit la sévérité et non la fréquence. Deux biais qui
s'annulent sur l'espérance peuvent ne pas s'annuler sur le quantile, et c'est le quantile qui
fait le SCR.}

% =====================================================================
\part*{II. Combien d'incidents par an}
\addcontentsline{toc}{part}{II. Combien d'incidents par an}

\section{Le processus de Poisson, et pourquoi il ne suffit pas}

\lequoi{Modéliser un nombre d'événements par an, sans savoir quand ils tombent.}

\lidee{Si les événements arrivent indépendamment, à un rythme constant, et jamais deux
exactement en même temps, alors leur nombre annuel suit une loi de Poisson. C'est le modèle
de comptage par défaut, et il n'a qu'un paramètre.}

\laformule{$\mathbb{P}(N=k)=e^{-\lambda}\lambda^k/k!$, avec la propriété caractéristique
\[
  \mathbb{E}[N]=\mathrm{Var}(N)=\lambda .
\]
Moyenne et variance sont \emph{égales}. C'est à la fois sa force et sa faiblesse.}

\cheznous{Cette égalité est testable, et elle est rejetée massivement. Sur $14\,944$
cellules entité-année, un test du rapport de vraisemblance donne $p\approx10^{-30}$.
Autrement dit~: la probabilité d'observer nos données si Poisson était vrai est de l'ordre
d'un sur $10^{30}$.}

\laphrase{Poisson impose que la variance égale la moyenne. Sur les cellules firme-année, la
variance dépasse la moyenne d'un facteur $1{,}19$ : cela paraît peu, mais sur quinze mille
cellules c'est un rejet à $10^{-30}$. Donc j'ai changé de loi.}

\lepiege{Ne pas dire \og Poisson sous-estime le risque\fg{}. Poisson donne la \emph{bonne}
moyenne, il se trompe sur la \emph{dispersion}, donc sur la queue. Le SCR est un quantile
extrême~: c'est la dispersion qui le fait, pas la moyenne.}

\section{La surdispersion et la négative binomiale}

\lequoi{Les données montrent beaucoup plus de variabilité que Poisson n'en autorise. Pourquoi,
et comment en tenir compte ?}

\lidee{Parce que le rythme lui-même varie. Certaines années sont plus exposées que
d'autres, certaines entités plus que d'autres. On rend donc $\lambda$ \emph{aléatoire} au
lieu de le fixer. Un mélange de Poisson par une loi Gamma donne exactement la négative
binomiale.}

\laformule{Si $N\mid\Lambda\sim\mathcal{P}(\Lambda)$ et $\Lambda\sim\mathrm{Gamma}$, alors
$N$ est négative binomiale et
\[
  \mathbb{E}[N]=\lambda, \qquad \mathrm{Var}(N)=\varphi\,\lambda, \qquad \varphi>1 .
\]
Le facteur $\varphi$ est la \textbf{surdispersion}~: le rapport variance sur moyenne.}

\cheznous{L'indice de dispersion d'un mélange de Poisson vaut $1+\lambda/r$ : il n'est pas
invariant d'échelle. Il vaut $1{,}19$ sur les cellules firme-année, $2{,}24$ sur la série annuelle
détendancée, et le moteur retient $\varphi=9{,}20$, choix de prudence posé et non mesure. Ce
$9{,}20$ et la charge systémique $a=0{,}60$ du modèle multi-états décrivent la même surdispersion
agrégée à $5\,\%$ près ($a=0{,}57$). À variance égale, la \emph{forme} du mélange (Gamma ou
lognormale) déplace encore la queue du nombre de $11\,\%$ ($74$ contre $82$).}

\laphrase{La négative binomiale, c'est Poisson dont le rythme est lui-même incertain.
Concrètement, certaines années sont plus exposées que d'autres, et ce sont ces années-là
qui font le capital.}

\lepiege{Le changement de loi de comptage \textbf{ne déplace pas la moyenne}, il déplace la
VaR, et de $+9\,\%$ seulement : à $\xi\approx0{,}6$, le quantile annuel est porté par la perte
unique dominante. Sur la série détendancée, la surdispersion mesurée ($2{,}24$) est quatre fois
plus faible que celle retenue ($9{,}20$) : le calage est conservateur, pas calibré, et il faut le
dire soi-même.}

\section{Pourquoi pas un processus de Hawkes}

\lequoi{Un processus auto-excité, où chaque événement augmente temporairement la
probabilité du suivant, serait un candidat naturel pour de la contagion.}

\lidee{Hawkes modélise la contagion \emph{dans le temps}~: un incident aujourd'hui rend un
autre incident plus probable demain. Notre cascade modélise la contagion \emph{entre
domaines}~: un incident sur un pilier rend un incident plus probable sur un autre pilier.
Ce ne sont pas les mêmes axes.}

\cheznous{Le test a été fait, sur la chronologie PRC datée, contre les noyaux de la
littérature (Bessy-Roland, Boumezoued et Hillairet ; Boumezoued, Cherkaoui et Hillairet).
L'excitation mesurée est intra-journalière, avec un pic à $3{,}3$~heures, et le ratio de
branchement de $0{,}551$ tombe vers zéro dès que les co-occurrences du même jour sont traitées
comme exogènes. Surtout, une cascade \emph{sans aucune} auto-excitation, simulée puis ajustée comme
une chronologie réelle, donne déjà $0{,}480$ au jour : à la résolution disponible, les deux
modèles ne sont pas distinguables.}

\laphrase{Hawkes décrit la contagion dans le temps, ma cascade décrit la contagion entre domaines.
Je l'ai testé contre les noyaux de la littérature, et à la résolution des données il n'est pas
distinguable d'une cascade sans auto-excitation. Je retiens la cascade par parcimonie, parce
qu'elle porte en plus la structure entre piliers.}

\lepiege{Ne pas dire que Hawkes est rejeté ou réfuté : c'est une équivalence observationnelle
suivie d'un choix de parcimonie. Attention au sens : dégrader la résolution fait \emph{monter} le
ratio de branchement estimé. Et ne pas présenter ce choix comme une critique de cette littérature,
qui travaille sur des données horodatées plus finement.}

% =====================================================================
\part*{III. Combien coûte un incident}
\addcontentsline{toc}{part}{III. Combien coûte un incident}

\section{Pourquoi la théorie des valeurs extrêmes}

\lequoi{Le SCR est un quantile à $99{,}5\,\%$. On veut donc décrire précisément la région
où l'on a le moins d'observations. Ajuster une loi sur l'ensemble des données optimise
l'ajustement là où il y a le plus de points, c'est-à-dire au \emph{centre}, exactement là
où l'on s'en moque.}

\lidee{On abandonne l'idée d'une loi unique pour tout. On modélise \emph{seulement} la
queue, avec une famille dont on peut démontrer qu'elle est la seule limite possible.}

\laphrase{Une loi ajustée sur toutes les données est calée sur le centre, or je n'ai besoin
que de la queue. La théorie des valeurs extrêmes dit quelle famille de lois est la seule
possible tout au bout, et je n'ajuste que celle-là.}

\section{Les deux théorèmes fondateurs}

\lidee{Deux résultats, structurés exactement comme le théorème central limite~: ils disent
qu'\emph{une seule} famille peut apparaître à la limite, quelle que soit la loi de départ.}

\laformule{\textbf{Fisher-Tippett-Gnedenko.} Si le maximum renormalisé de $n$ variables
indépendantes converge en loi, la limite ne peut être que la loi des valeurs extrêmes
généralisée~:
\[
  G_\xi(x)=\exp\!\left\{-\left(1+\xi x\right)^{-1/\xi}\right\}, \qquad 1+\xi x>0 .
\]
\textbf{Pickands-Balkema-de Haan.} De façon équivalente, la loi des \emph{excès} au-dessus
d'un seuil $u$ converge, quand $u$ monte vers le point terminal, vers la loi de Pareto
généralisée~:
\[
  \mathbb{P}(X-u\le y \mid X>u)\;\longrightarrow\;
  1-\left(1+\frac{\xi y}{\sigma}\right)^{-1/\xi}.
\]}

\cheznous{C'est le second qu'on utilise, parce qu'il exploite \emph{tous} les dépassements
de seuil et pas seulement le maximum de chaque bloc. On perd beaucoup moins d'information.}

\laphrase{C'est le même genre de résultat que le théorème central limite. Là où le TCL dit
que les moyennes tendent vers une gaussienne quelle que soit la loi de départ, ces deux
théorèmes disent que les extrêmes tendent vers une seule famille. Je n'ai donc pas
\emph{choisi} la loi de queue, elle est imposée.}

\lepiege{Ces théorèmes sont \emph{asymptotiques}. Ils ne garantissent pas que la GPD est
correcte à un seuil fini~: ils garantissent qu'aucune autre famille ne peut l'être à la
limite. C'est plus faible, il faut le dire, et c'est pourquoi on teste l'adéquation.}

\section{La méthode des excès et la loi de Pareto généralisée}

\lidee{On fixe un seuil $u$. En dessous, on garde la distribution empirique telle quelle. Au
dessus, on ajuste une GPD. La loi obtenue est dite \emph{recollée} (spliced).}

\laformule{La fonction de survie s'écrit, pour $x>u$~:
\[
  \mathbb{P}(X>x)\;=\;p_u\left(1+\frac{\xi (x-u)}{\sigma}\right)^{-1/\xi},
  \qquad p_u=\mathbb{P}(X>u).
\]}

\cheznous{$\xi=0{,}5954$ sur OpRisk, $\xi=1{,}03$ sur PRC. Le paramètre $\xi$ est
l'\textbf{indice de queue}, et il commande tout~:
\begin{itemize}[itemsep=1pt,topsep=2pt]
\item $\xi<1/2$~: variance finie ;
\item $1/2\le\xi<1$~: variance \emph{infinie}, moyenne finie : c'est le cas de la calibration
publiée ;
\item $\xi\ge1$~: moyenne elle-même infinie.
\end{itemize}
Le $\xi=1{,}03$ de la PRC dépasse $1$, mais il ne mesure pas un risque plus lourd : la sévérité
PRC est une puissance du nombre d'enregistrements, et une puissance multiplie l'indice de queue par
son exposant. C'est pourtant ce régime qui fait tomber la TVaR comme mesure de couverture : sous
lui, elle n'existe pas.}

\laphrase{L'indice de queue mesure à quel point les catastrophes dominent. Sur ma calibration il
vaut $0{,}6$ : la variance des pertes est infinie, leur moyenne finie, et un seul sinistre peut
peser plus que tous les autres réunis. C'est ce qu'on observe en cyber : sur OpRisk, le décile
supérieur porte plus de $92\,\%$ de la perte totale.}

\lepiege{Avec $\xi$ proche de $1$, la moyenne empirique n'est pas un estimateur stable et
les intervalles de confiance sont larges par nature. Ne jamais annoncer un SCR ponctuel
sans sa bande dans ce régime.}

\section{Le choix du seuil}

\lequoi{Le seuil $u$ arbitre entre biais et variance~: trop bas, on inclut le corps de la
distribution et l'approximation GPD est fausse ; trop haut, il reste trop peu de points et
l'estimation devient bruitée.}

\lidee{On ne le choisit pas au hasard : on le lit sur des diagnostics, puis on vérifie qu'une
règle écrite le retrouve.}

\laformule{\textbf{Vie résiduelle moyenne.} Pour une GPD, la fonction
$e(u)=\mathbb{E}[X-u\mid X>u]$ est \emph{linéaire} en $u$ : on choisit $u$ au début de la zone
linéaire. \textbf{Diagnostic GPD.} On trace $\hat\xi$ en fonction de $u$ : le seuil doit se
situer là où $\xi$ passe sous $1$ (espérance finie) en gardant au moins trente excès.}

\cheznous{Seuil retenu $u=20{,}03$~M\euro{} sur le périmètre cyber $\times$ finance d'OpRisk, soit
$91$ excès, et $p_u=0{,}1509$, valeur gelée (le recomptage donnerait $91/582=0{,}1564$, écart
chiffré au chapitre~11). Une règle de stabilité écrite, appliquée à dix percentiles candidats,
retrouve $19{,}87$~M\euro{} avec le même quantile ; la règle la plus permissive descendrait à
$12{,}88$ et donnerait un quantile supérieur de $16\,\%$. Attention : $\hat\xi$ n'est stable qu'au
\emph{voisinage} du seuil ; de $0{,}60$ il descend à $0{,}38$ quand on remonte au percentile $95$,
dans l'intervalle déjà publié.}

\cheznous{\textbf{Complément après le dépôt (script~119, hors rapport).} Une règle de stabilité
choisit un seuil ; un test d'adéquation le contrôle. J'ai calibré le test de
Kolmogorov-Smirnov : la table usuelle est trop indulgente quand les paramètres sont estimés sur
les données testées. La $p$-value vient donc d'un bootstrap paramétrique, où la GPD est
réestimée sur chaque échantillon simulé ($1\,000$ tirages par seuil, dix seuils de $8{,}62$ à
$105{,}39$~M\euro{}). Au seuil retenu ($91$ excès), $D=0{,}043$, $p=0{,}938$ pour
Kolmogorov-Smirnov et $p=0{,}963$ pour Anderson-Darling. Seuls les deux seuils les plus bas
($8{,}62$ et $10{,}95$) sont rejetés à $5\,\%$, et par Anderson-Darling seulement. Au-dessus de
$20$~M\euro{}, $\hat\xi$ reste dans l'intervalle à $95\,\%$ de chaque seuil. Les quatre estimateurs
au seuil retenu, avec leur intervalle à $95\,\%$ par bootstrap de $2\,000$ tirages des excès :
maximum de vraisemblance $0{,}598$ $[0{,}250\,;0{,}890]$, moments pondérés $0{,}507$
$[0{,}255\,;0{,}622]$, moments de Dekkers, Einmahl et de Haan $0{,}864$ $[0{,}524\,;1{,}081]$,
Hill $1{,}377$ $[1{,}183\,;1{,}585]$. Le quantile de sévérité à $99{,}5\,\%$, au taux gelé,
vaut $642$ $[369\,;1\,091]$ par maximum de vraisemblance et $560$ $[360\,;774]$ par moments
pondérés. La méthode reprend la structure du mémoire de F.~Dountio (test et bootstrap), pas ses
valeurs : sa variable est un nombre de dossiers, la nôtre une perte en euros.}

\lepiege{Une $p$-value élevée à haut seuil peut venir du faible nombre d'excès, donc d'une
puissance faible : le test écarte des seuils, il ne désigne pas un optimum. La sélection du
seuil sur les mêmes données n'est pas corrigée. Et le réajustement sur la base courante donne
$\hat\xi=0{,}598$ et un quantile de $642$, contre $0{,}595$ et $663$ publiés : les paramètres
publiés viennent de l'échantillon 2000-2026 ($86$ excès).}

\laphrase{Le seuil se lit sur la vie résiduelle moyenne, qui devient droite au-delà de vingt
millions, et une règle de stabilité écrite le retrouve à un pour cent près. Et il ne m'arrange pas :
la règle la plus permissive donnerait un capital plus élevé.}

\lepiege{Ne pas dire que le graphique de Hill confirme le seuil : Hill vaut $1{,}32$ au seuil
retenu, plus du double du maximum de vraisemblance, et le mémoire en fait une corroboration d'un
autre genre, puisque c'est exactement ce que Hill produit en échantillon fini sous $\xi=0{,}595$.
Et ne pas confondre avec les $4{,}176$~M\euro{} d'une table descriptive : c'est le percentile $85$
de la sévérité PRC dérivée, jamais un seuil.}

% =====================================================================
\part*{IV. Comment les pertes s'additionnent}
\addcontentsline{toc}{part}{IV. Comment les pertes s'additionnent}

\section{Le modèle collectif}

\lidee{On ne modélise pas la perte annuelle directement. On modélise séparément
\emph{combien} d'incidents et \emph{combien coûte} chacun, puis on somme. C'est le cadre
standard en actuariat non-vie.}

\laformule{
\[
  L \;=\; \sum_{i=1}^{N} X_i,
\]
avec $N$ le nombre d'incidents (fréquence) et $X_i$ leurs coûts (sévérité), les $X_i$
indépendants de $N$ et entre eux.}

\cheznous{Des dizaines de milliers d'années simulées par Monte-Carlo, avec le bruit de simulation
publié à côté de chaque grandeur. La loi de $L$ n'a pas de forme fermée simple avec une queue
lourde et une cascade, mais elle se calcule sans simuler par inversion de Fourier de sa fonction
génératrice : c'est le contrôle du moteur, qui retrouve le quantile à moins d'une erreur type
($5\,945$ contre $6\,049$ au conforme).}

\laphrase{Je sépare deux questions qui n'ont rien à voir~: combien d'incidents par an, et
combien coûte un incident. Chacune se calibre sur ses propres données, et la perte annuelle
est leur somme.}

\lepiege{L'indépendance entre $N$ et les $X_i$ est une hypothèse, et elle a été testée hors
échantillon : elle tient sur les excès. La pente négative qu'on voit sur toute la distribution
disparaît au-delà du seuil, c'est un artefact de profondeur de collecte. Et la cascade ne rend pas
$N$ et $X$ dépendants : elle fait dépendre le coût d'un sinistre du nombre de piliers qu'il
touche.}

\section{VaR, TVaR, et pourquoi le SCR est une VaR}

\lidee{La VaR à $99{,}5\,\%$ est le montant qu'on ne dépasse que dans une année sur deux
cents. La TVaR est la moyenne des pertes \emph{au-delà} de ce seuil.}

\laformule{
\[
  \mathrm{VaR}_\alpha(L)=\inf\{x:\ \mathbb{P}(L\le x)\ge\alpha\},
  \qquad
  \mathrm{TVaR}_\alpha(L)=\mathbb{E}[L\mid L>\mathrm{VaR}_\alpha(L)].
\]}

\cheznous{Solvabilité~II impose la VaR à $99{,}5\,\%$ à horizon un an, et le mémoire la prend
\emph{non centrée} : la perte attendue d'une non-conformité n'étant couverte par aucune provision,
elle n'est pas déduite. La TVaR sert de diagnostic, jamais de couverture : noyau de l'allocation
d'Euler, et indicateur de forme, $\mathrm{TVaR}/\mathrm{VaR}\approx2{,}1$ contre l'asymptote
$1/(1-\xi)=2{,}5$.}

\laphrase{Le capital, c'est la perte annuelle qu'on ne dépasse qu'une année sur deux cents. Ce
n'est pas mon choix de mesure, c'est celui de Solvabilité~II. La TVaR, je m'en sers pour
diagnostiquer la queue, pas pour couvrir.}

\lepiege{La VaR n'est pas sous-additive, la TVaR est cohérente : c'est vrai, et le mémoire ne dit
pas le contraire. Mais l'argument décisif contre une couverture en TVaR n'est pas de degré, il est
d'existence : sous la calibration PRC, $\xi=1{,}03>1$, l'espérance est infinie et la TVaR n'existe
pas. Une mesure de couverture qui cesse d'être définie selon la source ne peut pas porter un
capital.}

% =====================================================================
\part*{V. Comment un problème en entraîne un autre}
\addcontentsline{toc}{part}{V. Comment un problème en entraîne un autre}

\section{La convention d'indices, à ne jamais improviser}

\lequoi{Une matrice dirigée a deux lectures possibles, et les deux sont défendables. Si le
document en utilise l'une ici et l'autre là, tous ses énoncés directionnels deviennent
ambigus.}

\lidee{On fixe la convention une fois, on l'écrit, et on la vérifie sur un fait connu.}

\laformule{\[
  W_{jk}\;=\;\text{propension de } j \text{ (\textbf{source}) à entraîner la défaillance de }
  k \text{ (\textbf{cible}).}
\]
La \textbf{ligne} $j$ décrit ce que $j$ \emph{émet}, la \textbf{colonne} $k$ ce que $k$
\emph{reçoit}. La somme de ligne $s_j$ mesure donc une \textbf{émission}.}

\cheznous{Vérification sur la matrice d'expert : P1 a la plus grande somme de ligne ($2{,}6$)
et une somme de colonne trois fois inférieure à celle de P2. C'est exactement la revendication
centrale \og{}P1 émet fort et reçoit peu\fg{}, et elle serait \emph{fausse} sous la convention
inverse. La forme réduite du modèle latent s'écrit alors $(I-W^{\top})^{-1}$, le rayon spectral
étant de toute façon invariant par transposition.}

\laphrase{La ligne émet, la colonne reçoit. Je le vérifie sur ma matrice : le pilier que je dis
émetteur a bien la plus grosse somme de ligne. Une convention n'est pas un détail de notation,
c'est ce qui rend un énoncé directionnel falsifiable.}

\lepiege{Ce défaut a réellement existé dans le mémoire : l'équation structurelle sommait sur
les lignes alors que la définition de la cascade et tout le code sommaient sur les colonnes. Le
code et le classeur ont tranché. Vérifier une convention sur un fait connu coûte trente
secondes et évite une matrice transposée.}

\section{La matrice de contagion et sa normalisation}

\lequoi{On veut dire \og un problème sur le pilier $j$ aggrave le pilier $k$\fg{}, avec une
intensité. Mais si on met des coefficients au hasard, la cascade peut exploser.}

\begin{center}
\begin{tikzpicture}[scale=1.0]
\node[pilier] (p1) at (90:2.0)  {P1};
\node[pilier] (p2) at (18:2.0)  {P2};
\node[pilier] (p3) at (-54:2.0) {P3};
\node[pilier] (p4) at (234:2.0) {P4};
\node[pilier] (p5) at (162:2.0) {P5};
\node[etiq,above=1pt of p1] {gouvernance TIC};
\node[etiq,above right=0pt and 1pt of p2] {incidents};
\node[etiq,below=1pt of p3] {tests};
\node[etiq,below=1pt of p4] {tiers};
\node[etiq,left=1pt of p5]  {partage};
% P1 emet fort, recoit peu
\draw[flechef] (p1) to[bend left=12] (p2);
\draw[flechef] (p1) to[bend left=12] (p3);
\draw[flechef] (p1) to[bend right=12] (p4);
\draw[fleche]  (p4) to[bend right=12] (p2);
\draw[fleche]  (p4) to[bend left=12]  (p3);
\draw[flechep] (p5) to[bend left=12]  (p1);
\draw[flechep] (p3) to[bend left=12]  (p2);
\node[align=left,font=\scriptsize,text=navy!75,anchor=west] at (4.3,0.9)
  {\textbf{Lire la matrice~:}\\$W_{jk}$ = propension de $j$\\à entraîner $k$.\\
   \emph{La ligne émet, la colonne reçoit.}};
\node[align=left,font=\scriptsize,text=accent,anchor=west] at (4.3,-0.9)
  {P1 émet fort et reçoit peu.\\C'est pourquoi il sort\\premier en vue Shapley.};
\end{tikzpicture}
\end{center}

\noindent La cascade est un \emph{processus de branchement}~: un incident amorcé sur un pilier
tire, pour chaque arc, une pièce de Bernoulli de probabilité $W_{jk}$, et recommence depuis les
piliers atteints. Sans garde-fou, ce processus peut ne jamais s'éteindre.

\lidee{On part du jugement dirigé (la matrice \textsc{trans} du classeur), et on la
renormalise pour garantir mathématiquement que la propagation s'éteint.}

\laformule{
\[
  W \;=\; g\,\frac{T}{\max_j \sum_k T_{jk}},
  \qquad g=0{,}90 ,
\]
ce qui garantit $\rho(W)\le g<1$, où $\rho$ est le rayon spectral (la plus grande valeur
propre en module).}

\cheznous{Le facteur $g$ est un curseur de propagation~: $g=0$ signifie aucune contagion,
$g\to1$ signifie une cascade au bord de l'explosion. On le fixe à $0{,}90$, ce qui est
volontairement élevé.}

\laphrase{La normalisation garantit que la cascade s'éteint. Sans elle, un choc pourrait
s'amplifier indéfiniment en tournant en boucle entre les piliers, et le capital serait
infini.}

\lepiege{Les valeurs de $g$, $0{,}45$ au conforme, $0{,}68$ au partiel et $0{,}90$ au non
conforme, sont des \emph{choix}, pas des mesures. À dire ainsi : \og elles sont posées, mais la
signe est conservé sur les cinq correspondances testées et sur la grille en $g$, avec la cascade
indépendante. Cette vérification ne démontre pas la monotonie pour toute valeur de $g$, ni pour
tous les moteurs. L'amplitude reste un scénario\fg{}.}

\section{L'inverse de Leontief et la sous-criticité}

\lidee{Le choc initial se propage, ce qu'il déclenche se propage à son tour, et ainsi de
suite. La somme de tous les tours est une série géométrique matricielle.}

\laformule{
\[
  (I-W)^{-1} \;=\; I + W + W^2 + W^3 + \cdots
\]
Cette série converge \emph{si et seulement si} $\rho(W)<1$. Le terme $I$ est le choc
initial, $W$ ce qu'il déclenche directement, $W^2$ ce que cela déclenche à son tour.}

\cheznous{L'objet vient de l'économie (Leontief, tableaux entrées-sorties)~: la même algèbre
décrit comment une commande dans un secteur en génère dans les autres. Sur la matrice du mémoire,
$\rho(W)=0{,}562\,g$, soit $0{,}506$ à $g=0{,}9$ ; le seuil critique $\rho(W)=1$ serait atteint à
$g=1{,}78$, hors du domaine admissible.}

\laphrase{C'est la même formule que les tableaux entrées-sorties en économie. Le premier
terme est le choc direct, le deuxième ce qu'il déclenche, le troisième ce que cela déclenche
à son tour. La condition de convergence est exactement la condition pour que la cascade
s'éteigne.}

\lepiege{Le passage à $W^2$ est celui où la \emph{direction} commence à compter, alors
qu'elle ne compte pas dans $W$ seul si on ne regarde que les co-occurrences. C'est le point
technique le plus important du mémoire, voir la section~\ref{sec:trace}.}

\section{Le seuil de déclenchement, et le rôle de la loi normale}

\lequoi{À partir de quand un pilier bascule ? Il faut un mécanisme, pas seulement une
probabilité.}

\lidee{On suppose une variable latente $X_j$ qui mesure la fragilité du pilier $j$, et le
pilier bascule quand elle franchit un seuil. C'est le mécanisme de Merton, repris par
Vasicek pour le risque de crédit.}

\laformule{Le pilier $j$ bascule quand
\[
  X_j \;\ge\; K_j \;=\; F^{-1}(1-p_j),
\]
où $p_j$ est sa probabilité marginale de basculement et $F$ la fonction de répartition de la
latente.}

\cheznous{$K_4=1{,}03$ pour le pilier des prestataires, dont la probabilité de déclenchement
vaut $p_4=0{,}151$.}

\laphrase{Il y a deux questions derrière le mot seuil, et il ne faut pas les mélanger. Pour
savoir \emph{quand un pilier bascule tout seul}, la loi normale est inoffensive~: elle ne
sert que d'unité de mesure, la vraie queue lourde est dans la sévérité, pas ici. Pour savoir
\emph{combien de piliers basculent ensemble}, c'est autre chose, et là le gaussien pose un
vrai problème.}

\cheznous{\textbf{Un seuil par entité ?} Le cadre général est $K_{ij}=F_{ij}^{-1}(1-p_{ij})$ : $i$
est l'entité, $j$ le pilier, $F_{ij}$ la loi du stress latent et $p_{ij}$ la probabilité cible de
dépassement. Écrire $F_i$ seul suppose la même marge pour les cinq piliers d'une entité, ce que la
propagation ne garantit pas. Avec des chocs gaussiens, $\sigma_{ij}^2=[(I-W^{\top})^{-1}\Sigma_i
(I-W)^{-1}]_{jj}$ et $K_{ij}=\sigma_{ij}\Phi^{-1}(1-p_{ij})$. L'écriture est possible ;
l'estimation ne l'est pas : aucune base n'observe la perte par pilier ni par entité, et un écart de
$0{,}10$ sur une sensibilité demande environ $813$ périodes. Le mémoire garde donc un seuil $K_j$
commun, et le capital passe par la cascade, une sévérité commune puis la descente d'échelle.}

\lepiege{À probabilités et dépendance inchangées, transformer ensemble une marge et son seuil ne
change aucun déclenchement. Changer $F_i$ ne personnalise donc rien : il faudrait des $p_{ij}$
différents, une autre dépendance, ou une autre loi de perte. Et cette loi latente, sans unité, ne
remplace pas la GPD des sévérités en euros.}

\lepiege{Le choix de la loi de la latente est une \emph{normalisation} pour la marge~:
n'importe quelle loi continue donnerait la même probabilité marginale, seul le seuil change
de valeur. Ce n'est pas vrai pour la dépendance. Confondre les deux est l'erreur classique.}

\section{Les copules et la dépendance de queue}

\lequoi{Comment plusieurs piliers basculent-ils \emph{ensemble} ? La probabilité marginale
de chacun ne suffit pas à le dire.}

\lidee{Une copule sépare les marges de la structure de dépendance. On peut garder les mêmes
marges et changer complètement la façon dont les événements extrêmes se coordonnent.}

\laformule{Le coefficient de \textbf{dépendance de queue supérieure} mesure la
co-occurrence extrême~:
\[
  \lambda_U \;=\; \lim_{q\to1^-}\ \mathbb{P}\bigl(U_2>q \mid U_1>q\bigr).
\]
Pour la copule gaussienne, $\lambda_U = 0$ dès que la corrélation est $<1$. Pour la copule
de Gumbel de paramètre $\theta$,
\[
  \lambda_U \;=\; 2-2^{1/\theta}.
\]}

\cheznous{Avec $\theta=1{,}8$, on a $\lambda_U = 0{,}53$. À $\tau$ de Kendall \emph{égal},
la gaussienne donne $0$ et Gumbel $0{,}53$. La probabilité que les cinq piliers basculent
ensemble passe de $0{,}17$ à $0{,}39$ selon la copule retenue. On encadre au lieu de
trancher~: plancher gaussien, plafond de Gumbel. Le co-déclenchement d'un autre pilier avec P4
tient alors dans $[0{,}48\,;0{,}58]$, sous la borne $\gamma=0{,}68$ du choc commun, qui reçoit
ainsi son origine structurelle.}

\laphrase{Une dépendance gaussienne a une propriété gênante~: dans les extrêmes, elle prédit
que les choses se \emph{découplent}. C'est le contraire de ce qu'on observe en crise. Donc
je n'ai pas tranché entre les deux, j'ai encadré, et l'écart entre le plancher et le plafond
est un résultat en soi.}

\lepiege{Deux copules peuvent avoir le \emph{même} $\tau$ de Kendall, donc la même
dépendance \og moyenne\fg{}, et des comportements extrêmes opposés. Comparer deux copules par
leur corrélation ne dit rien sur la queue.}

\section{Les états de conformité : Vasicek à l'échelle rapide, Markov à l'échelle lente}
\label{sec:conformite}

\lequoi{Être conforme à DORA n'est ni un oui ni un non, et cela change avec le temps. Comment
mettre un état de conformité dans un modèle de capital, sans le figer et sans le réduire à un
interrupteur ?}

\lidee{On sépare deux horloges. À l'échelle \emph{rapide}, celle des incidents et de leur
cascade dans l'année, l'état de conformité est une photographie : chaque pilier est non
conforme (NC), partiellement conforme (PC) ou conforme (C), et le moteur calcule le capital
\emph{sachant} cette configuration. Cet état est tiré par un Vasicek. À l'échelle \emph{lente},
celle des mois et des années, une chaîne de Markov fait progresser la conformité, et son état
règle les paramètres de la couche rapide. On sépare les deux pour qu'une tendance longue
n'explique pas un enchaînement de quelques heures.}

\begin{center}
\begin{tikzpicture}[node distance=7mm]
\node[brique] (nc) {NC\\non conforme};
\node[brique,right=of nc] (pc) {PC\\partiellement};
\node[brique,right=of pc] (c) {C\\conforme};
\draw[fleche] (nc)--(pc); \draw[fleche] (pc)--(c);
\node[etiq,above=2mm of pc] {échelle lente : chaîne de Markov, séjours d'Erlang-$k$, C absorbant};
\node[brique,below=17mm of nc] (v) {Vasicek\\multi-états};
\node[brique,right=of v] (q) {Quatre canaux\\$\lambda,\ g,\ p_u$, accumulation};
\node[brique,right=of q,fill=vert!15,draw=vert!70!black] (k) {Capital\\sachant l'état};
\draw[fleche] (v)--(q); \draw[fleche] (q)--(k);
\draw[flechef] (pc.south) -- node[etiq,right,pos=0.45] {l'état règle les canaux} (q.north);
\node[etiq,below=2mm of q] {échelle rapide : les incidents et leur cascade dans l'année};
\end{tikzpicture}
\end{center}

\laformule{\textbf{D'où vient le seuil.} Merton pose qu'un défaut survient quand une variable
latente franchit un seuil. Vasicek y ajoute un facteur commun, qui crée la dépendance :
\[
  X_i \;=\; \sqrt{\rho}\,Y + \sqrt{1-\rho}\,\varepsilon_i,
  \qquad \text{défaut si } X_i \ge K_i = \Phi^{-1}(1-\mathrm{PD}_i).
\]
\textbf{Notre version, à trois états.} Chaque pilier a une capacité de conformité latente,
décomposée à la Vasicek en un facteur systémique $\Theta$ et un choc propre :
\[
  C^\star \;=\; \gamma\,\Theta + \sqrt{1-\gamma^2}\,\varepsilon,
  \qquad \gamma = 0{,}68 .
\]
Au lieu d'un seuil de défaut, on en pose deux, ordonnés, $K_{\text{bas}} < K_{\text{haut}}$ :
\[
  \text{NC si } C^\star \le K_{\text{bas}}, \qquad
  \text{PC si } K_{\text{bas}} < C^\star \le K_{\text{haut}}, \qquad
  \text{C si } C^\star > K_{\text{haut}},
\]
avec $K_{\text{bas}} = \Phi^{-1}(p_{NC})$ et $K_{\text{haut}} = \Phi^{-1}(p_{NC}+p_{PC})$. C'est
un probit ordonné. Sachant l'environnement $\Theta$ :
\[
  \mathbb{P}(\mathrm{NC}\mid\Theta) = \Phi\!\Bigl(\tfrac{K_{\text{bas}}-\gamma\Theta}{\sqrt{1-\gamma^2}}\Bigr),
  \qquad
  \mathbb{P}(\mathrm{C}\mid\Theta) = 1-\Phi\!\Bigl(\tfrac{K_{\text{haut}}-\gamma\Theta}{\sqrt{1-\gamma^2}}\Bigr).
\]}

\cheznous{Les probabilités d'état sont ancrées sur les enquêtes de conformité de fin 2025 :
NC $35\,\%$, PC $35\,\%$, C $30\,\%$. Les cinq piliers partagent $\Theta$ avec la charge
$\gamma=0{,}68$, soit une corrélation $\gamma^2 = 0{,}462$ entre deux piliers quelconques. En
crise systémique ($\Theta$ très négatif), la masse bascule vers la non-conformité. L'état agit
ensuite par quatre canaux : la fréquence $\lambda$ passe de $21{,}56$ à $53{,}62$ incidents par an,
le gain de propagation $g$ de $0{,}45$ (C) à $0{,}68$ (PC) puis $0{,}90$ (NC), le taux de
matérialité $p_u$ est multiplié par $0{,}85$ (C) ou $1{,}20$ (NC), et l'accumulation par un
prestataire commun apparaît à l'état NC, avec une charge de $0{,}68$.}

\laformule{\textbf{L'échelle lente.} Une chaîne de Markov à temps continu NC $\to$ PC $\to$ C,
dont l'état conforme est absorbant dans le cas de base. Les séjours sont de type phase
(Erlang-$k$) : une exponentielle autoriserait une remédiation quasi instantanée, alors qu'un
projet DORA a une durée minimale. Avec $Q$ le générateur et $v_0$ la distribution initiale,
\[
  \bigl(\mathbb{P}(\mathrm{NC},t),\,\mathbb{P}(\mathrm{PC},t),\,\mathbb{P}(\mathrm{C},t)\bigr)
  \;=\; v_0\, e^{Qt}\ \text{(sommes de phases)},
  \qquad
  \mathrm{SCR}(t) \;=\; \sum_{e\in\{\mathrm{NC},\mathrm{PC},\mathrm{C}\}} \mathbb{P}(e,t)\,\mathrm{SCR}_e .
\]
\textbf{Le couplage.} Le même $\Theta$ agit sur les deux couches : un environnement dégradé
ralentit la remédiation, $\mu(\Theta)=\mu_0\,e^{\beta_{\text{rem}}\Theta}$, et durcit la cascade,
$\mathrm{SCR}_e(\Theta)=\mathrm{SCR}_e\,e^{-\beta_{\text{casc}}\Theta}$. La dispersion de $\Theta$
engendre la bande d'incertitude de la trajectoire.}

\cheznous{De l'état marginal actuel (NC $35\,\%$, PC $35\,\%$, C $30\,\%$), le capital OpRisk
passe de $10\,752$ à $7\,041$~M\euro{} en cinq ans. La formule $\mathrm{SCR}(t)$ est une moyenne
de quantiles : elle diffère du quantile à 99,5\,\% de la loi mélangée, de $9{,}3\,\%$ sur le
mélange des $32$ configurations. La trajectoire se lit donc comme l'espérance du besoin de capital
sous l'incertitude d'état.}

\laphrase{À l'échelle rapide, l'état de conformité est un probit ordonné à la Vasicek : une
capacité latente et deux seuils qui découpent trois états. À l'échelle lente, une chaîne de
Markov le fait progresser de NC à C, et son état règle les quatre paramètres du moteur. Le seuil
vient de Merton, le facteur commun de Vasicek, la migration d'état du risque de crédit : la
chaîne joue le rôle de la migration d'un rating de conformité, le Vasicek celui du défaut
conditionnel au rating.}

\lepiege{\textbf{Deux latentes de sens opposé.} Dans la cascade, $X$ mesure le \emph{stress} :
un $X$ élevé est une mauvaise nouvelle. Ici $C^\star$ est une \emph{capacité} : élevée, elle
est bonne, et la non-conformité est en bas. Le facteur $\Theta$ de conformité n'est pas non plus
le facteur $Y_j$ par pilier de la cascade.}

\lepiege{\textbf{Les trois valeurs de $g$ sont posées}, parce que l'état de conformité DORA des
entités n'est pas publié. La cascade indépendante respecte l'ordre du capital sur la grille de vingt
points et les cinq correspondances testées. Cela ne démontre ni une monotonie pour tout $g$ ni la
même propriété pour le moteur de marche. L'écart \og $116$ à $169$~M\euro{}\fg{} reste conditionnel
aux valeurs posées.}

\lepiege{La loi d'Erlang est \emph{posée}. Elle n'a aucun rapport avec la loi mélangeante de la
fréquence du chapitre sur le modèle fréquence-sévérité, dont la forme estimée est $0{,}63$. Et
les données ne mesurent pas le niveau de conformité d'une entité : c'est pourquoi l'état est
supposé, jamais observé.}

% =====================================================================
\part*{VI. Ce qu'on ne peut pas savoir, et pourquoi}
\addcontentsline{toc}{part}{VI. Ce qu'on ne peut pas savoir, et pourquoi}

\section{La décomposition $W = S + A$}

\lidee{Toute matrice se coupe de manière unique en une partie symétrique et une partie
antisymétrique. La première dit \emph{qui est lié à qui}, la seconde dit \emph{dans quel
sens}.}

\laformule{
\[
  S=\tfrac12\bigl(W+W^{\top}\bigr), \qquad A=\tfrac12\bigl(W-W^{\top}\bigr),
  \qquad W=S+A .
\]
Les deux parties sont \textbf{orthogonales}~: $\langle S,A\rangle=0$, d'où un Pythagore
$\|W\|^2=\|S\|^2+\|A\|^2$.}

\cheznous{Sur la matrice ordonnée estimée à partir du corpus de rapports, la part de direction
$\|A\|/\|W\|$ vaut $60\,\%$, et elle ne bouge pas quand on corrige le dénominateur, alors que le
niveau des $p_{jk}$ tombe de $0{,}27$ à $0{,}07$. La direction n'est pas un détail de la matrice.}

\laphrase{On coupe la matrice en deux~: une partie qui dit quels piliers sont liés, et une partie
qui dit dans quel sens. Les deux sont orthogonales, et la co-occurrence que voit la donnée ne
dépend, au premier ordre, que de la première.}

\lepiege{La décomposition propre porte sur le \emph{flux stationnaire} $\pi_j P_{jk}$, pas
sur la matrice de transition brute. Notre $S=(W+W^{\top})/2$ en est le cas où $\pi$ est
uniforme, et l'écart mesuré est de $0{,}069$, donc faible. Le dire soi-même met à l'abri.}

\medskip
\noindent\textbf{Le schéma qui fait comprendre la thèse en dix secondes.} Deux matrices
opposées, une seule statistique observable, deux décisions contraires.

\begin{center}
\begin{tikzpicture}
% cas W
\node[pilier] (a1) at (0,1.1) {P1};
\node[pilier] (b1) at (0,-0.5) {P2};
\draw[flechef] (a1) -- (b1);
\node[font=\scriptsize\bfseries,text=navy] at (0,1.9) {matrice $W$};
\node[etiq,align=center] at (0,-1.3) {P1 entraîne P2\\\emph{remédier P1}};
% cas W transpose
\node[pilier] (a2) at (3.4,1.1) {P1};
\node[pilier] (b2) at (3.4,-0.5) {P2};
\draw[flechef] (b2) -- (a2);
\node[font=\scriptsize\bfseries,text=navy] at (3.4,1.9) {matrice $W^{\top}$};
\node[etiq,align=center] at (3.4,-1.3) {P2 entraîne P1\\\emph{remédier P2}};
% ce que la donnee voit
\node[draw=vert!70!black,fill=vert!8,rounded corners=2pt,align=center,
      font=\footnotesize,inner sep=6pt] (obs) at (8.4,0.3)
  {\textbf{Ce que la donnée agrégée voit}\\
   \og P1 et P2 tombent ensemble\fg{}\\[2pt]
   au premier ordre, $\Prob(j \text{ et } k)$\\
   ne dépend que de $S$, et $S(W)=S(W^{\top})$};
\draw[fleche] (1.0,0.3) -- (obs.west |- 0,0.3);
\draw[fleche] (4.4,0.3) -- (obs.west |- 0,0.3);
\node[align=center,font=\scriptsize,text=accent] at (4.25,-2.3)
  {\textbf{Indiscernables pour la donnée. Matériellement différentes pour la décision.}};
\end{tikzpicture}
\end{center}

\noindent C'est l'énoncé de la frontière d'identification, et il ne repose sur \emph{aucun} dire
d'expert : au premier ordre et sous amorce uniforme, la co-occurrence est une fonction de $S$ seule,
et $S$ est identique pour $W$ et pour sa transposée. \textcolor{accent}{Ce n'est pas une
impossibilité de principe} : au-delà du premier ordre, la loi complète des piliers touchés par
incident distingue $W$ de $W^\top$, une fourche d'un collisionneur ($62{,}31\,\%$ de sinistres
multi-piliers sous $W$, $50{,}77\,\%$ sous $W^\top$). La frontière est celle des statistiques que
les sources fournissent, et c'est pourquoi le registre recommandé consigne tous les piliers
touchés.

\section{Le renversement du temps et la fluctuation martingale}

\lequoi{Pourquoi n'arrive-t-on pas à mesurer $A$ ? Et est-ce une limite de méthode ou une
limite de principe ?}

\lidee{Passer le film à l'envers. La partie symétrique ne change pas, la partie
antisymétrique s'inverse. Et la fluctuation d'un processus, une fois la tendance retranchée,
ne dépend que de la première.}

\laformule{La chaîne renversée a pour flux $F^{*}=F^{\top}=S-A$. Et pour toute fonction test
$f$, l'intensité de la fluctuation vaut
\[
  \mathbb{E}_\pi\bigl[\Gamma(f,f)\bigr]
  \;=\;\sum_{j,k} S_{jk}\bigl(f(k)-f(j)\bigr)^2 ,
\]
la partie $A$ y contribuant exactement zéro. La démonstration tient en une ligne~:
$(f(k)-f(j))^2$ est symétrique en $(j,k)$, $A$ est antisymétrique, les termes $(j,k)$ et
$(k,j)$ se compensent.}

\cheznous{Vérifié sur trois cents fonctions test tirées au hasard : l'énergie de fluctuation est
la même pour la chaîne posée et pour sa symétrisée, les points tombent sur la diagonale (figure du
chapitre~8, \S~8.3).}

\laphrase{Renverser le temps laisse la partie symétrique intacte et retourne le signe de la partie
directionnelle. Or la fluctuation d'un processus, sa partie martingale, ne dépend que de la
symétrique. Donc une statistique invariante par renversement du temps ne peut pas voir la
direction. La loi complète des piliers touchés par incident, elle, le pourrait : c'est une donnée
qui manque, pas une impossibilité.}

\lepiege{Ne pas dire que $W$ \emph{est} une chaîne de Markov. C'est un noyau de contagion~: on
emprunte l'algèbre de la réversibilité comme cadre de lecture, on n'ajoute pas une hypothèse. Et le
mot martingale s'entend au sens des chaînes de Markov, sous la mesure physique : aucune mesure
risque-neutre, aucun test de martingalité.}

\section{Le test placebo}
\label{sec:placebo}

\lequoi{Même des données sans aucune direction produisent un déséquilibre, par simple
fluctuation d'échantillonnage. Observer un déséquilibre ne prouve donc rien tant qu'on ne
sait pas ce que le hasard produit tout seul.}

\lidee{On fabrique une version des données dont la direction a été \textbf{détruite
volontairement}, tout le reste étant conservé, et on compare. C'est l'essai clinique~: si le
médicament ne fait pas mieux que la pilule vide, il ne fait rien.}

\laformule{On mesure l'asymétrie de la matrice de transitions comptées~:
\[
  \mathcal{A}(M)=\frac{\|M-M^{\top}\|_1}{2}
  =\frac12\sum_{j,k}\bigl|M_{jk}-M_{kj}\bigr| ,
\]
c'est-à-dire le surplus de transitions d'un sens sur l'autre, sommé sur toutes les paires.
Puis on \textbf{permute les années à l'intérieur de chaque entité}~: l'ensemble des
catégories touchées chaque année est conservé, seul l'ordre est détruit. On recompte, on
recalcule, et on compare~:
\[
  z=\frac{\mathcal{A}(M_{\text{obs}})-\overline{\mathcal{A}^{\,\text{null}}}}
        {\widehat{s}(\mathcal{A}^{\,\text{null}})} .
\]}

\cheznous{
\begin{center}\small
\begin{tabular}{@{}l r r r l@{}}
\toprule
\textbf{Source} & \textbf{observée} & \textbf{sous placebo} & \textbf{$z$} & \textbf{Verdict}\\
\midrule
OpRisk, pas annuel & $119$ & $128\pm27$ & $-0{,}33$ & aucun signal \\
7 post-mortems & \multicolumn{2}{c}{\emph{même protocole}} & $+3{,}93$ & signal net \\
\bottomrule
\end{tabular}
\end{center}
Sur OpRisk, l'observé est \emph{sous} la moyenne du bruit. Un test binomial par paire le
confirme~: aucune paire ne survit à la correction de Bonferroni.}

\laphrase{Je voulais savoir si la direction que je crois voir est réelle. Le problème, c'est
que même des données au hasard produisent un déséquilibre. Alors j'ai pris mes données, j'ai
détruit l'ordre en gardant tout le reste, et j'ai comparé. Mes vraies données font
$119$, le hasard fait $128$~: elles font moins bien que le hasard. Puis j'ai refait
exactement le même test sur les sept rapports d'enquête, et là le signal apparaît nettement.
Même test, autre source, résultat opposé~: le problème venait de la donnée, pas de la
méthode.}

\lepiege{Trois choses à ne pas dire. \textbf{(i)} Le test du \og temps inversé\fg{} est
\emph{dégénéré} sur un comptage~: renverser la chronologie donne $M^{\top}$ par
construction, donc une corrélation de $-1$ quel que soit le contenu. Il ne prouve rien.
\textbf{(ii)} Un test non rejeté n'est pas une preuve d'absence~: la donnée est
\emph{compatible} avec l'absence de direction. \textbf{(iii)} Sur les post-mortems, le test
établit la \emph{cohérence} du codage, pas son \emph{exogénéité}~: le biais de narration des
rapports d'enquête n'est pas levé par le calcul.}

\section{La production d'entropie}

\lidee{Un seul nombre qui dit à quel point le processus est dirigé, et qui a une
interprétation exacte~: à quel point le film à l'endroit est distinguable du film à
l'envers.}

\laformule{
\[
  \sigma=\sum_{j<k} 2A_{jk}\,\ln\!\frac{S_{jk}+A_{jk}}{S_{jk}-A_{jk}} \;\ge\; 0,
\]
nulle si et seulement si $A=0$. Elle est aussi le taux d'entropie relative entre la loi de
la trajectoire lue à l'endroit et celle lue à l'envers.}

\cheznous{$\sigma$ est strictement positive pour la chaîne posée et s'annule à la
symétrisation (figure du chapitre~8, \S~8.3). \textbf{Et la conséquence qui compte}, qui se lit
sur la formule : chaque terme de la somme est \emph{pair} en $A_{jk}$, donc $\sigma$ ne dépend que
des \emph{amplitudes} de la direction, pas de ses signes. Les $1\,024$ sommets de l'ensemble
admissible, qui ne diffèrent que par les signes, ont donc la même production d'entropie, alors
que le capital y varie d'un facteur $1{,}27$.}

\laphrase{La production d'entropie mesure à quel point on peut distinguer le film à l'endroit du
film à l'envers. Si elle est nulle, personne ne peut dire dans quel sens le temps coule. Mais
attention~: elle ne voit que la \emph{quantité} de direction, jamais son \emph{sens}. Les mille
vingt-quatre configurations ont la même production d'entropie, et pourtant elles étalent le
capital de plus d'un quart.}

\lepiege{Conséquence à tirer~: ce n'est pas $\sigma$ qu'il faut chercher à estimer, ce sont
les \emph{signes}, un par un. C'est ce qui justifie que la valeur d'information se
hiérarchise par paire et non globalement.}

\section{$z$ scalaire ou $z$ par paire}

\lequoi{L'écriture $P=S+zA$ est ambiguë tant qu'on n'a pas dit si $z$ est un nombre ou dix.}

\laformule{\textbf{Lecture globale}~: $z\in[-1,1]$ règle l'\emph{intensité} d'un motif de
direction fixé. Un seul degré de liberté. Elle ne peut pas changer quel pilier précède
quel autre.
\textbf{Lecture par paire}, celle du mémoire~:
\[
  A_{jk}=z_{jk}\,t\,S_{jk}, \qquad z_{jk}\in[-1,1], \qquad j<k,
\]
soit \textbf{dix degrés de liberté}. Le paramètre $t$ est l'amplitude maximale d'asymétrie,
les $z_{jk}$ portent chacun un sens indépendant.}

\cheznous{Les $2^{10}=1024$ sommets sont exactement les configurations $z_{jk}=\pm1$. La
contrainte de positivité du flux impose $|A_{jk}|\le S_{jk}$, donc $t\le1$~; $t=1$ est
l'ignorance totale.}

\laphrase{Un $z$ global ne fait que renforcer ou affaiblir un ordre déjà choisi. Ce qu'il me
faut, c'est un curseur par paire, parce que c'est le \emph{sens} de chaque lien que je ne
connais pas, pas son intensité globale.}

\section{Le troisième échec : le schéma existe, il n'est pas rempli}

\lequoi{Deux bases ont échoué à identifier la direction, l'une faute de taxonomie, l'autre
faute de résolution temporelle. On pouvait encore croire que personne n'avait \emph{prévu} de
collecter la bonne donnée.}

\lidee{On teste l'hypothèse la plus favorable : un standard de description d'incidents qui,
lui, prévoit exactement les bons champs. Le schéma \textsc{veris}, référence de la profession
et support du \emph{Data Breach Investigations Report}, porte littéralement les deux objets
manquants : un champ \texttt{control\_failure} qui nomme le domaine de contrôle défaillant, et
une chronologie en quatre étapes qui porte l'antériorité.}

\cheznous{Sur les $10\,591$ incidents de la base publique associée : \texttt{control\_failure}
renseigné dans \textbf{$0{,}12\,\%$} des cas ($13$ incidents, dont un seul en secteur
financier), chronologie datée sur au moins deux étapes dans \textbf{$1{,}19\,\%$}. Et
l'\textbf{intersection} des deux, qui est ce qu'identifier $W$ exigerait simultanément, ne
compte qu'\textbf{un seul incident}, et \textbf{aucun} en secteur financier.}

\laphrase{Le champ qui donnerait la réponse existe dans le standard de la profession. Il est
rempli treize fois sur dix mille incidents. Le problème n'est donc pas de savoir quoi collecter, il
est d'en
rendre le remplissage obligatoire, et c'est précisément ce qu'un régime de notification
réglementaire peut imposer là où une base déclarative volontaire ne le peut pas.}

\lepiege{C'est le plus utile des trois échecs, parce qu'il transforme la recommandation. On ne
demande plus d'inventer une taxonomie, ce qui serait long et contestable, on demande de rendre
obligatoires deux champs \textbf{déjà normalisés}. Coût de conception nul.}

\section{La matrice ordonnée $p_{jk}$}

\lequoi{Le codage documentaire, tel qu'on l'exploite d'abord, compare $n(j\to k)$ à
$n(k\to j)$ \emph{à l'intérieur} d'une paire. Ce rapport élimine le dénominateur, donc il
donne un \emph{sens} et ne dit rien de l'\emph{amplitude}.}

\lidee{On ajoute le dénominateur. Chaque incident déclare non seulement ses transitions, mais
l'ensemble des \textbf{piliers défaillants}, y compris ceux qui ont cédé sans rien entraîner.
On peut alors estimer directement chaque case de la matrice ordonnée.}

\laformule{Avec $N_j$ le nombre d'incidents où $j$ a cédé et $m_{jk}$ ceux où le rapport
établit $j\to k$ :
\[
  p_{jk}=\Prob(k \text{ cède}\mid j \text{ a cédé}),
  \qquad
  p_{jk}\mid\text{données}\ \sim\ \mathrm{Beta}\bigl(1+m_{jk},\ 1+N_j-m_{jk}\bigr).
\]
C'est la matrice ordonnée \emph{complète}, donc $S$ \textbf{et} $A$.}

\cheznous{Sur dix rapports et $22$ transitions : \textbf{douze} liens sur vingt ont un crédible
à $90\,\%$ entièrement \emph{sous} $0{,}5$, et \textbf{aucun} au-dessus. Les plus élevés,
$P_3\to P_2$ ($0{,}77$, $[0{,}48\,;0{,}95]$) et $P_1\to P_4$ ($0{,}68$, $[0{,}40\,;0{,}89]$),
recouvrent encore $0{,}5$. Les douze liens faibles comprennent les \emph{quatre} sorties de
$P_2$, soit toutes, et trois des quatre entrées de $P_1$, l'arc $P_5\to P_1$ n'étant pas établi
comme faible : le puits est établi comme tel, la source pure seulement en partie.}

\laphrase{Mon corpus établit très bien où la contagion \emph{ne va pas}, et pas du tout où elle
va fort. Le puits et la source pure sont donc établis en tant que tels, mais aucun coefficient
n'est établi comme grand. C'est un ordre, ce n'est pas une magnitude.}

\lepiege{Une enquête raconte la \emph{chaîne} qui a produit le sinistre : elle documente bien
les contrôles qui ont cédé \emph{et} propagé, mal ceux qui ont cédé sans propager, pas du tout
ceux qui ont tenu. Le dénominateur $N_j$ est donc sous-estimé et $p_{jk}$ est une \textbf{borne
supérieure}. En ajoutant des défaillances non propagées, le \emph{niveau} s'effondre de
$0{,}27$ à $0{,}07$ tandis que la part de direction $\lVert A\rVert/\lVert W\rVert$ ne bouge
pas ($60\,\%$). Le corpus établit un ordre, pas une magnitude, et cette phrase-là est robuste.}

\section{L'identification partielle}

\lidee{Quand la donnée ne détermine pas un paramètre, on peut étudier les résultats compatibles
avec un ensemble d'hypothèses déclaré. Ici, la partie symétrique $S$ est posée à dire d'expert ;
seule la direction varie. L'exploration numérique est donc conditionnelle à ce choix.}

\laformule{On définit l'ensemble admissible
\[
  \mathcal{W}=\bigl\{\,W=S+A \ :\ |A_{jk}|\le t\,S_{jk}\ \bigr\},
\]
et l'on rapporte $\bigl[\inf_{\mathcal{W}}\mathrm{SCR},\ \sup_{\mathcal{W}}\mathrm{SCR}\bigr]$
au lieu d'un point.}

\cheznous{Socle sans contagion $5\,275$~M\euro. Avec contagion et ignorance totale, plage calculée aux sommets
dans $[6\,858\,;\,8\,697]$~M\euro, un facteur $1{,}27$ : la contagion ajoute entre $+30$ et
$+65\,\%$ au socle. L'énumération porte sur $1\,024$ sommets, et $3\,000$ points intérieurs
tirés au hasard ne sortent pas de l'enveloppe des sommets, ni en VaR ni en perte moyenne (bornée
dans $[849\,;1\,100]$). Le point d'expert, $8\,110$, est un point de l'ensemble, jamais sa
mesure.}

\laphrase{Je ne prétends pas connaître la direction, donc je ne donne pas un chiffre, je
donne une fourchette. Ce n'est pas de la prudence, c'est la seule réponse possible~: une
grandeur qui compte pour le résultat et qui est invisible dans mes données.}

\lepiege{L'atteinte des extrêmes aux sommets est une \emph{propriété vérifiée
numériquement}, pas un théorème démontré. C'est déclaré comme telle dans le mémoire.}

\medskip
\noindent\textbf{Le schéma à projeter si l'on vous demande \og quel est votre SCR\fg{}.}

\begin{center}
\begin{tikzpicture}[xscale=1.0]
\def\y{0}
% axe
\draw[->,draw=navy!60] (-0.4,\y-1.25) -- (11.2,\y-1.25) node[right,font=\scriptsize,text=navy]{M\euro};
\foreach \x/\v in {0/5275, 3.6/6858, 7.4/8697} {
  \draw[draw=navy!60] (\x,\y-1.15) -- (\x,\y-1.35);
  \node[font=\scriptsize,text=navy] at (\x,\y-1.62) {$\v$};
}
% socle
\fill[vert!25,draw=vert!70!black] (0,\y-0.35) rectangle (3.6,\y+0.35);
\node[font=\scriptsize,align=center,text=vert!35!black] at (1.8,\y) {\textbf{socle}\\sans contagion};
% bande d'ignorance
\fill[accent!18,draw=accent!80!black] (3.6,\y-0.35) rectangle (7.4,\y+0.35);
\node[font=\scriptsize,align=center,text=accent!70!black] at (5.5,\y)
  {\textbf{bande d'ignorance}\\largeur $1\,839$};
% resserrement
\draw[decorate,decoration={brace,amplitude=4pt},draw=navy] (7.4,\y+0.5) -- (5.1,\y+0.5)
  node[midway,above=4pt,font=\scriptsize,text=navy,align=center]{un registre spécifié\\en retire $66\,\%$};
\draw[densely dashed,draw=navy!70] (5.1,\y+0.36) -- (5.1,\y+0.5);
\node[align=left,font=\scriptsize,text=navy,anchor=west] at (7.9,\y+0.15)
  {$1\,024$ sommets\\énumérés};
\end{tikzpicture}
\end{center}

\noindent La largeur de la bande n'est pas une incertitude subie~: c'est une grandeur
\textbf{calculée à l'avance}, et le mémoire dit combien elle se réduirait si telle donnée
existait. C'est ce qui transforme une limite en livrable.

\section{Pourquoi $A$ compte quand même}
\label{sec:trace}

\lequoi{Si $A$ est invisible dans la donnée, pourquoi s'en soucier ?}

\laformule{Parce que le capital passe par $(I-W)^{-1}=I+W+W^2+\cdots$, et dès le terme
quadratique $W^2=S^2+SA+AS+A^2$ fait apparaître $A$. En prenant la trace~:
\[
  \operatorname{tr}(W^2)\;=\;\|S\|^2-\|A\|^2 .
\]
Le signe moins vient de l'antisymétrie~: $A^{\top}=-A$ implique $A^2=-AA^{\top}$.}

\laphrase{La direction est invisible au premier ordre mais elle apparaît dès qu'on met la
matrice au carré, et le calcul du capital passe par là. Une grandeur qui compte pour la
réponse et qui est absente de la donnée~: c'est la définition exacte d'une identification
partielle.}

% =====================================================================
\part*{VII. À quel point on y croit}
\addcontentsline{toc}{part}{VII. À quel point on y croit}

\section{Le bootstrap}

\lidee{On ne connaît pas la loi de l'estimateur. Alors on la fabrique~: on rééchantillonne
les données avec remise, on réestime à chaque fois, et la dispersion des estimations
obtenues approche la dispersion vraie.}

\cheznous{L'intervalle de confiance bootstrap du quantile de sévérité à $99{,}5\,\%$,
$[411\,;1\,037]$~M\euro, vaut un \textbf{facteur $2{,}5$} entre ses deux bouts. Le bootstrap est
aussi utilisé en version
\emph{paramétrique} pour les tests d'adéquation~: on simule sous la loi ajustée pour obtenir
la loi de la statistique de test.}

\laphrase{Le bootstrap fabrique la loi de mon estimateur en rééchantillonnant mes propres
données. Et le résultat est brutal~: rien que cette incertitude-là vaut un facteur deux et
demi sur le quantile d'un sinistre, et près de neuf sur la charge annuelle.}

\lepiege{Le bootstrap capture l'incertitude \emph{de paramètre} à modèle fixé. Il ne dit rien du
risque que le modèle lui-même soit faux. Et le facteur $2{,}5$ porte sur le quantile d'\emph{un}
sinistre : porté sur la charge annuelle, le même intervalle de $\xi$ donne un facteur $8{,}93$. Ne
jamais comparer tel quel ce $2{,}5$ à un facteur lu sur la charge annuelle.}

\section{Les tests d'adéquation}

\lidee{Une loi ajustée peut être fausse. On teste la distance entre la loi ajustée et les
données, et on regarde si cette distance est anormalement grande.}

\laformule{Kolmogorov-Smirnov mesure l'écart maximal entre les fonctions de répartition.
Anderson-Darling pondère ce même écart pour donner davantage de poids aux \emph{queues}, ce
qui est exactement la région qui nous intéresse.}

\cheznous{Par bootstrap paramétrique : Anderson-Darling $p=0{,}96$ et Kolmogorov-Smirnov
$p=0{,}89$ pour la famille GPD ; avec les paramètres publiés imposés, $0{,}99$ et $0{,}89$, ce
qui atteste le chiffre et pas seulement le choix de loi. La couverture réelle de l'intervalle à
$90\,\%$ de $\xi$ vaut $86\,\%$ : l'intervalle est \emph{trop étroit}, et il faudrait élargir sa
demi-largeur d'environ $10\,\%$. La binomiale négative bat Poisson avec $p\approx10^{-30}$.}

\laphrase{Les tests ne rejettent pas, avec des valeurs $p$ de $0{,}89$ à $0{,}99$. Attention à ce
que cela veut dire~: les \emph{formes} ne sont pas rejetées, mais c'est le \emph{niveau} qui reste une
bande, et la couverture des quantiles de sévérité est, elle, rejetée hors échantillon.}

\lepiege{Les valeurs critiques tabulées de KS et AD sont fausses quand les paramètres sont
\emph{estimés} sur les mêmes données. Il faut passer par un bootstrap paramétrique, ce que
nous faisons, et il faut le dire, parce que beaucoup de mémoires ne le font pas.}

\section{Les trois étages d'incertitude}

\lidee{L'incertitude ne se réduit pas à l'intervalle de confiance. Elle a trois sources qui
se cumulent et dont les tailles n'ont rien à voir.}

\cheznous{
\begin{center}\small
\begin{tabular}{@{}l l r r@{}}
\toprule
\textbf{Étage, sur la charge annuelle} & \textbf{Ce qui varie} & \textbf{Bande (M\euro)} &
\textbf{Facteur} \\
\midrule
Paramètre & $\xi$ sur son intervalle à $90\,\%$ & $[3\,042\,;27\,152]$ & $8{,}93$ \\
Identification & la direction, non identifiée & $[6\,858\,;8\,697]$ & $1{,}27$ \\
Modèle, dépendance & de l'indépendance au comonotone & $[5\,322\,;9\,806]$ & $1{,}84$ \\
Modèle, famille de queue & de la lognormale à la queue lourde & $[4\,330\,;39\,901]$ & $9{,}21$ \\
\bottomrule
\end{tabular}
\end{center}}

\laphrase{Lues sur la même grandeur, la charge annuelle, l'incertitude de paramètre et celle de
modèle sont du même ordre, un facteur de l'ordre de neuf chacune, et elles écrasent celle de
l'identification. La bande que je reporte cumule le paramètre et l'identification ; le choix de
famille, qui ne se moyenne pas, est publié comme axe de sensibilité. Et le résultat qui tient est
ailleurs : \textbf{le niveau bouge dans la bande, l'ordre d'amorce des piliers ne bouge pas d'une
famille à l'autre}.}

\lepiege{Deux fautes à éviter. Comparer le facteur $2{,}5$ du quantile d'un sinistre au facteur
$5{,}5$ de la famille de queue (lognormale $4\,330$ contre GPD à $\xi=0{,}90$, $24\,016$) : ce sont
deux échelles. Et laisser croire que rien ne tient : toujours enchaîner avec les tests d'adéquation
et l'invariance de l'ordre. Les largeurs de paramètre font varier $\xi$ seul : elles ne sont ni des intervalles conjoints
ni des majorants garantis de l'incertitude totale.}

\section{VaR prédictive et VaR robuste}

\lequoi{La remarque de Caroline Hillairet~: prendre un quantile d'un objet incertain est
mauvais. Comment y répondre sans abandonner la VaR, qui est imposée ?}

\lidee{Deux réponses différentes, à ne pas confondre.
\textbf{La VaR prédictive} intègre l'incertitude de paramètre dans la loi elle-même~: on
mélange les lois sur la loi a posteriori des paramètres, puis on prend le quantile du
mélange.
\textbf{La VaR robuste} prend le pire cas sur un ensemble d'ambiguïté~: on ne mélange pas,
on borne.}

\laformule{
\[
  \mathrm{VaR}^{\text{préd}}_\alpha = \mathrm{VaR}_\alpha\!
  \left(\int f(\cdot\mid\theta)\,\pi(\theta)\,d\theta\right),
  \qquad
  \mathrm{VaR}^{\text{rob}}_\alpha = \sup_{\theta\in\Theta}\ \mathrm{VaR}_\alpha
  \bigl(f(\cdot\mid\theta)\bigr).
\]}

\cheznous{Six postures, en M\euro{} : point $657$ ; prédictive $645\pm7$ (le $648$ d'un autre
calcul est le même nombre à la précision du bruit) ; robuste à $\beta=90\,\%$ $939\pm13$, à
$95\,\%$ $1\,037\pm9$, à $99\,\%$ $1\,247\pm24$ ; pire cas de modèle, $\xi$ poussé à $0{,}90$ :
$1\,331$. \textbf{La prédictive coïncide presque avec le point} à $99{,}5\,\%$, à deux fois le
bruit ; elle ne s'en écarte qu'en profondeur, $+4\,\%$ à $99{,}9\,\%$. Et la robuste à $95\,\%$
est exactement la borne haute de l'intervalle déjà publié, $1\,037$ : la même quantité sous deux
noms. Le mémoire reporte donc le point avec sa bande.}

\laphrase{Caroline a raison, et c'est mesurable. Ma réponse n'est pas d'abandonner la VaR,
qui est imposée, mais de rendre l'incertitude visible sur trois étages. Et j'ai un résultat
un peu contre-intuitif~: la VaR prédictive coïncide presque avec la VaR ponctuelle. Ce n'est
pas l'intégration qui manquait, c'est la largeur de bande qui est le vrai sujet.}

\lepiege{Ne pas présenter la prédictive comme \og la bonne réponse\fg{} qui remplacerait le
point. Elle règle un problème (l'incertitude de paramètre) et pas l'autre (l'ambiguïté de modèle).
C'est la robuste qui règle le second, et les deux ne se substituent pas. Et la posture la plus
prudente est la moins précise : $\pm24$ à $\beta=99\,\%$, quand les deux extrémités de l'échelle
n'ont aucun bruit.}

% =====================================================================
\part*{VIII. Les outils annexes}
\addcontentsline{toc}{part}{VIII. Les outils annexes}

\section{Le kappa de Cohen}

\lequoi{Deux codeurs lisent les mêmes récits. Comment mesurer leur accord, sachant que même
deux personnes répondant au hasard seraient parfois d'accord par chance ?}

\laformule{
\[
  \kappa = \frac{p_o - p_e}{1 - p_e},
\]
où $p_o$ est l'accord observé et $p_e$ l'accord attendu par pur hasard. $\kappa=1$ est un
accord parfait, $\kappa=0$ un accord au niveau du hasard.}

\cheznous{Le dispositif de second codage en aveugle est prêt : dix récits réduits aux faits,
anonymisés, et une grille de $100$ cellules dont $22$ sont directionnelles dans le codage de
référence, ce qui fixe la base d'accord (matériel complémentaire). Sans second jeu, le script ne
fabrique aucun coefficient. Par convention, un $\kappa$ au-delà de $0{,}6$ se lit comme un accord
substantiel ; en dessous, cela voudrait dire que la séquence des défaillances ne se lit pas de
façon reproductible dans les rapports d'enquête, ce qui serait un résultat en soi.}

\laphrase{Le kappa corrige l'accord du hasard : deux personnes qui répondraient au hasard seraient
parfois d'accord, et le kappa retranche cette part. Le second codage n'est pas encore reçu ; le
dispositif est prêt, et il ne produit rien sans lui.}

\section{La valeur d'information}

\lidee{Une recommandation de reporting ne vaut rien tant qu'elle n'est pas chiffrée. On la
chiffre en mesurant de combien la bande de capital rétrécit quand on acquiert
l'information.}

\laformule{Avec $L(\mathcal{V})=\sup_{\mathcal{V}}\mathrm{SCR}-\inf_{\mathcal{V}}\mathrm{SCR}$,
\[
  V(\mathcal{I}) = 1-\frac{L(\mathcal{W}_{\mathcal{I}})}{L(\mathcal{W})} \in[0,1].
\]
Cette quantité est \emph{monotone} (une information plus riche ne vaut pas moins) et
\emph{sans dimension} (donc comparable d'une source de sévérité à l'autre).}

\cheznous{Si l'on retenait le sens que les sept premiers rapports documentent sur six
dépendances, la largeur de la plage tomberait de $32{,}7\,\%$ en moyenne sur quatre graines, de
$22$ à $44\,\%$ selon la graine, sur une grille commune où chaque coefficient vaut $-S$, $0$ ou $S$.
Ce n'est pas un acquis : le mémoire n'en revendique que deux paires, et ce chiffre est un
indicateur de ce qu'apporterait une information de cette étendue. Le premier calcul annonçait
$66\,\%$ : il fixait les amplitudes à leur maximum, alors qu'un sens connu les laisse libres entre
zéro et ce maximum. La valeur se \textbf{hiérarchise} sans se départager : quatre dépendances, dont
gouvernance-incidents ($13{,}3\,\%$), lèvent chacune de $11{,}2$ à $13{,}3\,\%$ en moyenne, leur
ordre variant entre graines ; les cinq dernières n'en lèvent que de $0{,}7$ à $2{,}3\,\%$.}

\laphrase{J'ai chiffré ce que vaut, en capital, une donnée qu'on n'a pas. Une information de sens
sur six dépendances resserrerait la plage d'un tiers en moyenne. Et surtout, quatre dépendances
valent chacune beaucoup plus que les cinq autres : on sait où la donnée manque, sans pouvoir dire
laquelle d'entre elles vaut le plus.}

\lepiege{Même les dix sens connus ($57{,}9\,\%$ de la plage levée), l'amplitude resterait inconnue. L'identification partielle n'est pas un état provisoire que la donnée
finirait par lever, elle est structurelle.}

\section{Le bayésien conjugué}

\lidee{Plutôt que d'estimer un signe de direction par oui ou non, on lui attribue une
probabilité, qu'on met à jour avec les observations documentaires.}

\laformule{Modèle Beta-binomial~: si $\theta\sim\mathrm{Beta}(\alpha_0,\beta_0)$ est la
probabilité que la direction aille dans un sens, et qu'on observe $n_{lo}$ transitions dans
ce sens et $n_{hi}$ dans l'autre, alors
\[
  \theta \mid \text{données} \ \sim\ \mathrm{Beta}(\alpha_0+n_{lo},\ \beta_0+n_{hi}),
\]
et l'on pose $A = S\,(2\theta-1)$, de sorte que $\theta=1/2$ donne $A=0$.}

\cheznous{Le crédible à $90\,\%$ du capital est $81\,\%$ plus resserré que les bornes de pire
cas. Sous un prior sceptique, la médiane ne bouge presque pas : la conclusion vient des comptes
documentés. Mais seulement $2$~paires documentées sur $6$ ont un intervalle qui exclut $0{,}50$, et
les quatre paires que le corpus ne documente pas gardent un $\theta$ uniforme : le posterieur
reproduit de lui-même l'identification partielle.}

\laphrase{Le bayésien me permet de dire \og ce lien va probablement dans ce sens\fg{} au lieu de
trancher par oui ou non. Deux résultats~: le prior n'a presque aucune influence, ce qui est
rassurant, et seulement deux paires sur six sont réellement tranchées, ce qui est honnête. Je
publie les deux objets : la plage aux sommets, conditionnelle à $S$ posée, et le crédible, qui ajoute un prior sur la direction.}

\lepiege{Le bayésien ne crée pas d'information. Il exprime autrement la même incertitude, de
façon plus lisible et plus fine, mais il ne remplace pas la donnée manquante.}

% =====================================================================
\part*{IX. Du capital à la décision}
\addcontentsline{toc}{part}{IX. Du capital à la décision}

\section{À quelle échelle se lit un SCR}

\lequoi{Un capital de plusieurs milliards paraît absurde pour une compagnie d'assurance. C'est
la première objection qu'on vous fera, et une clause de style n'y répond pas.}

\lidee{La fréquence d'amorce et l'échelle des pertes changent avec la taille, puis le quantile annuel
est recalculé. Le calcul d'entité utilise la cascade indépendante, sans choc commun P4 ; le chiffre
de tête à quatre canaux utilise la marche auto-évitante avec ce choc. Les montants issus de ces
deux protocoles ne se déduisent pas directement les uns des autres.}

\cheznous{\;\\[-6pt]
\begin{center}\small
\begin{tabular}{@{}l r r r r@{}}
\toprule
\textbf{Lecture} & $\lambda$/an & \textbf{Sévérité} & \textbf{SCR} & \textbf{Moyenne} \\
\midrule
Secteur financier mondial & $21{,}56$ & $\times1$ & $8\,123$~M\euro & $983$~M\euro \\
Seau \og{}au moins dix événements\fg{} & $0{,}210$ & $\times1$ & $488$~M\euro & $9{,}7$~M\euro \\
Seau \og{}toutes firmes\fg{} & $0{,}100$ & $\times1$ & $229$~M\euro & $4{,}6$~M\euro \\
$\lambda$ lu à la taille de l'entité & $0{,}092$ & $\times1$ & $198$~M\euro & $4{,}3$~M\euro \\
\textbf{Les deux canaux composés} & $0{,}092$ & $\times0{,}854$ & $\mathbf{169}$~M\euro & $3{,}6$~M\euro \\
\bottomrule
\end{tabular}
\end{center}
La dernière ligne est la lecture cohérente : fréquence \emph{et} sévérité prises à la taille de
l'entité. Sa bande sur les $1\,024$ sommets va de $131{,}5$ à $183{,}3$~M\euro.}

\laphrase{Les montants sectoriels concernent le périmètre financier étudié. Le calcul d'entité
donne $169$~M\euro{} dans la lecture du registre, sous un état supposé. C'est une extrapolation
d'ordre de grandeur : la forme de la sévérité reste celle de grandes institutions. Rien ne
garantit que ce montant majore le besoin réel. Le critère conventionnel de poids dans le SCR
ne définit pas un domaine de validité statistique.}

\lepiege{La non-linéarité est le vrai résultat. Diviser $\lambda$ par $235$ divise la moyenne par
$270$, donc à proportion, mais ne divise le capital que par $48$ : c'est le \textbf{principe du
grand saut unique}, le quantile d'une entité décrivant un sinistre isolé et non une accumulation.
Diviser le capital du secteur par le nombre d'entités le sous-estimerait d'un facteur de l'ordre de
cinq. Et le repère de $3\,\%$ des provisions n'est qu'un ordre de grandeur : il s'écarte de
$-8\,\%$ et $+56\,\%$ des modules publiés dans les SFCR.}

\lepiege{\textbf{La descente d'échelle ne distingue ni la taille fine ni le métier.} Les élasticités
à la taille valent $0{,}074$ pour la fréquence et $0{,}087$ pour la sévérité : une petite entité
subit presque la sévérité d'un grand groupe, donc sa charge rapportée au SCR est plus forte. Un
essai exploratoire, hors rapport, l'illustre. À taille égale, l'assurance a environ $0{,}5$ fois la
fréquence de la banque et environ $2{,}4$ fois sa sévérité ; un plafond par sinistre proportionnel à
l'exposition ($\kappa\times$ actifs) ferait suivre la charge à la taille. Mais le plafond $\kappa$
serait posé, la base ne couvre pas les petites entités, et le quantile à $99{,}5\,\%$ d'un petit
assureur tombe à zéro. Piste d'amélioration, pas résultat.}

\section{Détenir ou transférer}

\lequoi{Le mémoire compare son SCR à la Formule Standard, aux pertes réelles et au marché
cyber. Trois comparaisons de \emph{niveau}. Il en manque une de \emph{prix}.}

\lidee{Un SCR n'est pas une perte, c'est du capital \textbf{immobilisé}, et le capital coûte.
Solvabilité~II chiffre ce coût explicitement. En face, transférer coûte une prime. Les deux
grandeurs sont homogènes, donc comparables, et c'est l'arbitrage que fait un directeur des
risques.}

\laformule{Le coût annuel de détention vaut $\mathrm{CoC}\times\mathrm{SCR}$, où
$\mathrm{CoC}$ est le taux de coût du capital de la marge de risque : $6\,\%$ historiquement,
ramené à $\mathbf{4{,}75\,\%}$ par la directive (UE)~2025/2. Pour un traité en excès de perte
annuel de portée $L$, on minimise
\[
  \mathrm{co\hat ut}(L)=\underbrace{\E[\min(X,L)]/0{,}485}_{\text{prime}}
  \;+\;\mathrm{CoC}\times\underbrace{\mathrm{VaR}_{99{,}5\%}\bigl(X-\min(X,L)\bigr)}_{\text{capital résiduel}} .
\]}

\cheznous{Multiple de capital $\mathrm{SCR}/\E[X]=\mathbf{47}$ à l'échelle entité contre $8{,}3$
au secteur. Coût de détention $8{,}0$~M\euro/an au taux révisé ($10{,}1$ à $6\,\%$) contre
$3{,}6$~M\euro{} de sinistralité attendue, soit $2{,}2$ à $2{,}8$ fois. Le capital résiduel vaut
$\max(\mathrm{SCR}-L,0)$, donc l'optimum tombe exactement en $L^\star=\mathrm{SCR}=169$~M\euro, et
le coût de financement passe à $\mathbf{2{,}9}$~M\euro, soit $\mathbf{-64\,\%}$ ; pertes
conservées comprises, de $11{,}67$ à $5{,}12$~M\euro, soit $-56\,\%$. Une entité devenue conforme
(fréquence de $0{,}092$ à $0{,}046$) voit en outre ce coût optimal baisser de $60\,\%$ : la
conformité déplace le point d'équilibre entre rétention et transfert.}

\laphrase{Immobiliser le capital coûte deux fois et demie ce que le risque coûte en moyenne.
Près de l'optimum, un euro de portée coûte un peu plus d'un pour cent de prime et libère
quatre virgule sept cinq pour cent de coût du capital. C'est pour ça que le transfert domine, et
c'est un calcul, pas une opinion.}

\lepiege{Le \og coût total\fg{} s'appelle désormais coût de \emph{financement} (prime et coût du
capital), et ne comprend pas les pertes conservées. Ce gain est une \textbf{borne supérieure}, pour trois raisons : la capacité du marché
plafonne près de $500$~M\euro{} sur un seul risque ; les polices excluent la guerre et les
infrastructures essentielles ; et sous Solvabilité~II céder ne supprime pas l'exigence, cela
\emph{substitue} au risque cyber un risque de \textbf{défaut de contrepartie}. En résumant les
trois par la part de la portée achetée qui reste sans allègement de capital, le transfert reste
préférable, face au meilleur traité, tant que cette part est inférieure à
$\mathbf{64\,\%}$ ($81{,}5\,\%$ pertes conservées comprises) : il faudrait que près des deux tiers de
la couverture soient inefficaces pour renverser la conclusion. La conclusion est robuste, le
chiffre ne l'est pas.}

% =====================================================================

\input{cours_complements_institut_2026_10_28}

\end{document}
